% generated by figures/make_figures.py — do not edit
\begin{longtable}{lp{14.5cm}}
\toprule
Table & What it shows \\
\midrule
\endhead
\multicolumn{2}{l}{\emph{Pre-specified analyses}} \\
A1 & The secondary comparisons of the protocol under each split. \\
A2 & Ratio-bootstrap intervals of $R$ for every comparison and split. \\
A3 & Operating characteristics of the decision rule on synthetic labels for each split. \\
A4 & The primary comparison within descriptive strata (H5). \\
A5 & Secondary metrics of the held-out predictions. \\
A6 & The confirmatory comparisons with the frozen code and with the corrected reading of Norwegian names. \\
A7 & The coefficients of the second stage. \\
A8 & Interval coverage in the null simulations, the $p$ values of the protocol and the spatial correlation of residuals. \\
\multicolumn{2}{l}{\emph{Exploratory checks of the verdict}} \\
A9 & The corrected equivalence test on the ratio at four levels of resampling, its operating characteristics, and other measures of the change. \\
A10 & The equivalence test on the ratio, in the order the intervals were defined, with percentile and normal limits, TOST $p$, upper limits and the odds ratios they exclude. \\
A11 & Jackknife, studentized and weighted versions of the test, the effective number of clusters and the operating characteristics with activations resampled within countries. \\
A12 & The names increment over the recalibrated baseline model, decomposed into the regional composition and the names within areas. \\
A13 & The comparison on observable ground and the language-model reading with country indicators. \\
A14 & The names increment by AOI, continent and flood type. \\
A15 & The direction of the exploratory estimates that exclude the no-effect value. \\
A16 & The models of the exploratory analyses. \\
A17 & The within-area comparison under exploratory variants of sample, labels and reading. \\
A18 & The primary comparison in subsets of cells and under other fold assignments. \\
A19 & The achieved equivalence bound, the location groups needed to settle the inconclusive comparisons, the permutation test, the interval with the second stage refitted and the operating characteristics at the margin. \\
A20 & The Murphy decomposition, the comparison within AOIs and the skill of flood maps. \\
A21 & The primary comparison with one definition of permanent water, other subsampling seeds and a bare-earth elevation model. \\
A22 & The primary comparison by sensor, product generation and without Africa, and with the alternative second stages and references of amendment 10. \\
A23 & Brier scores of every split. \\
A24 & The activation split in two code versions, further intervals, the GFM domain and the denominator of the labels. \\
A25 & A census of the exploratory interval estimates. \\
A26 & Operating characteristics of the two decision rules, the reference model, constant forecasts, outlying null simulations, inclusion, flood typing, GFM images, built-up land near settlements, where the flooded cells near water-related names lie and a penalty on the names alone. \\
A27 & What the regional composition of names adds. \\
A28 & Whether areas with more water-related names flood more, within language regions. \\
A29 & Exploratory runs and controls of amendment 6. \\
\multicolumn{2}{l}{\emph{What assignment through settlements can carry}} \\
A30 & The injected odds ratio that each interval's upper limit excludes, under every no-effect reference. \\
A31 & The same gains with every model refitted on the injected labels. \\
A32 & How an effect on the ground a name refers to is diluted on all ground nearest the name. \\
A33 & Translations of the margin and bounds under the synthetic labels of the protocol. \\
A34 & Recurrent surface water around settlements, by name class (no flood labels). \\
\multicolumn{2}{l}{\emph{Flooding near settlements}} \\
A35 & Further checks of the odds ratios near settlements: fractional cover, flood type, population and built-up land, observable ground, a jackknife and the positional join. \\
A36 & Diagnosis of the GFM labels: flooded shares and baseline gains by AOI, and agreement with EMS by distance and built-up share. \\
A37--A38 & Where settlements with names of each class stand relative to the drainage, by distance (no flood labels). \\
A39 & Heights of settlements above the main rivers, by name class (no flood labels). \\
A40 & Within-AOI odds ratios of flooding by class and distance. \\
A41 & Odds ratios near settlements by class and distance band, with pairs-bootstrap intervals. \\
A42 & The names increment by class of the nearest name. \\
A43 & Flood contrasts and increments by distance from the settlement. \\
A44 & The pooled odds ratio near settlements under every variant, by river size and by class, and the two label sets on the same cells. \\
A45 & Where the flooded cells near water-related names lie, by country and by activation. \\
A46 & The pooled odds ratio near settlements with each activation left out. \\
A47 & Adjusted $p$ and $q$ values of the class coefficients. \\
A48 & The odds ratio near water-related names corrected over a grid of misreading rates. \\
A49 & The odds ratio near water-related names on every evaluable cell, with both label sets. \\
A50 & The floodplain reading by country, without one activation and without doubtful forms. \\
A51 & Every quantity of the analyses on GFM labels (amendment 3). \\
A52 & Diagnostics of the GFM labels. \\
A53 & Calibration and skill of the held-out predictions, by split. \\
\multicolumn{2}{l}{\emph{How the names were read}} \\
A54 & Homographs and wrong-language readings, the edited lexicon, Unicode names, the two coders and the language-model reading by block. \\
A55 & Elevation and river size of settlement pairs that differ by a relative qualifier. \\
A56 & The lexicon by language block. \\
A57 & The water-related forms that matched the most cells. \\
A58 & Forms excluded by the lexicon's own rules or prone to transferred names, and the most frequent forms by country. \\
A59 & Examples of names read as water-related, with their recorded derivations. \\
A60 & Precision and recall of the readings in more detail. \\
A61 & Recorded etymologies by lexicon class and language. \\
A62 & Agreement of a second coder and the precision and recall from its codes. \\
A63 & The classes of the language-model reading. \\
A64 & The settlements by GeoNames feature code, and the pilot simulations. \\
A65 & Settlements whose namesake is recorded on Wikidata. \\
A66 & Precision and recall by language. \\
A67 & Precision and recall by country and by form, for both coders. \\
A68 & Cells whose names are read in a language other than their own. \\
\multicolumn{2}{l}{\emph{Checks after the eleventh critiques (amendment 15)}} \\
A69 & The ground within 500 m of settlements with and without water-related names. \\
A70 & Census of the exploratory interval estimates of amendments 2 to 15. \\
\bottomrule
\end{longtable}
